\documentclass[aspectratio=169,10pt]{beamer}

\usetheme{Madrid}
\usecolortheme{default}

\usepackage{amsmath}
\usepackage{amssymb}
\usepackage{booktabs}
\usepackage{graphicx}

\definecolor{pgoBlue}{RGB}{37,84,150}
\definecolor{pgoGreen}{RGB}{54,130,91}
\definecolor{pgoOrange}{RGB}{191,104,43}

\setbeamercolor{structure}{fg=pgoBlue}
\setbeamercolor{alerted text}{fg=pgoOrange}
\setbeamertemplate{navigation symbols}{}
\setbeamertemplate{footline}{%
  \leavevmode%
  \hbox{%
    \begin{beamercolorbox}[
      wd=0.32\paperwidth,ht=2.5ex,dp=1.1ex,center
    ]{author in head/foot}%
      \usebeamerfont{author in head/foot}\insertshortauthor
    \end{beamercolorbox}%
    \begin{beamercolorbox}[
      wd=0.48\paperwidth,ht=2.5ex,dp=1.1ex,center
    ]{title in head/foot}%
      \usebeamerfont{title in head/foot}\insertsectionhead
    \end{beamercolorbox}%
    \begin{beamercolorbox}[
      wd=0.20\paperwidth,ht=2.5ex,dp=1.1ex,center
    ]{date in head/foot}%
      \usebeamerfont{date in head/foot}%
      \insertframenumber{} / \inserttotalframenumber
    \end{beamercolorbox}%
  }%
  \vskip0pt%
}

\newcommand{\norm}[1]{\left\lVert #1\right\rVert}
\newcommand{\eref}{E_{\mathrm{ref}}}
\newcommand{\Kff}{K_{\!ff}}
\newcommand{\Eup}{E_{up}}
\newcommand{\Eupf}{E_{up,f}}
\newcommand{\good}[1]{\textcolor{pgoGreen}{#1}}
\newcommand{\warn}[1]{\textcolor{pgoOrange}{#1}}

\title[Systematic Poking Calibration]{Systematic Poking Material Calibration}
\subtitle{From Direct Stress Data to Equilibrium Reaction Forces}
\author{libpgo material calibration experiments}
\date{August 2026}

\AtBeginSection[]{%
  \begin{frame}[plain]
    \centering
    \vfill
    {\usebeamercolor[fg]{structure}\Large
      Section \insertsectionnumber\ of 4}

    \vspace{1.2em}
    \makebox[\linewidth][c]{%
      \begin{beamercolorbox}[
        wd=0.84\paperwidth,sep=1.5em,center
      ]{title}
        {\LARGE\bfseries\insertsectionhead}
      \end{beamercolorbox}%
    }

    \vspace{1.4em}
    {\large From Direct Stress Data to Equilibrium Reaction Forces}
    \vfill
  \end{frame}
}

\begin{document}

\begin{frame}
  \titlepage
\end{frame}

\begin{frame}{One Story, Two Calibration Experiments}
  \small
  Both experiments fit the same 18-parameter Systematic Poking material to
  the same synthetic compressible Neo-Hookean target.

  \vspace{0.6em}
  \begin{columns}[T,onlytextwidth]
    \begin{column}{0.47\linewidth}
      \begin{block}{Experiment I: direct constitutive data}
        \[
          (F,P^\star)
          \longrightarrow
          P(F;p)
        \]
        The deformation $F$ is observed. This isolates material
        parameterization and stress sensitivities.
      \end{block}
    \end{column}
    \begin{column}{0.47\linewidth}
      \begin{block}{Experiment II: reaction-force data}
        \[
          (u_c,y^\star)
          \longrightarrow
          u^\star(p)
          \longrightarrow
          y(p)
        \]
        The interior state is unknown and must satisfy FEM equilibrium.
      \end{block}
    \end{column}
  \end{columns}

  \vspace{0.6em}
  \begin{alertblock}{Main question}
    Does the verified constitutive parameter pipeline remain correct and
    identifiable after introducing nonlinear static equilibrium?
  \end{alertblock}
\end{frame}

\begin{frame}{Roadmap}
  \tableofcontents[hideallsubsections]
\end{frame}

\section{Material Model and Shared Parameterization}

\begin{frame}{Target and Fitted Material}
  \small
  For principal stretches $s_1,s_2,s_3>0$ and $J=s_1s_2s_3$, the target is
  compressible Neo-Hookean:
  \[
    \psi_{\mathrm{NH}}
    =
    \sum_{i=1}^{3}
    \left[
      \frac{\mu}{2}(s_i^2-1)-\mu\log s_i
    \right]
    +\frac{\lambda_{\mathrm{NH}}}{2}(\log J)^2.
  \]
  The synthetic ground truth is
  $E^\star=2\times10^5$ and $\nu^\star=0.35$.

  \vspace{0.5em}
  The fitted paper-style Systematic Poking model is
  \[
    \boxed{
      \psi_{\mathrm{SP}}(s;p)
      =\sum_{i=1}^{3}f(s_i;q)+\lambda\,\widehat h(J)
    }.
  \]
  Its stretch curvature is a piecewise-linear spline,
  \[
    f''(s;q)=\sum_{k=0}^{16}q_k\phi_k(s),
    \qquad f(1)=f'(1)=0,
  \]
  while $\widehat h$ is a fixed log-squared volume spline shape.
\end{frame}

\begin{frame}{The 18 Optimizable Parameters}
  \small
  The physical parameter vector is
  \[
    p=
    \begin{bmatrix}q_0&\cdots&q_{16}&\lambda\end{bmatrix}^{T}
    \in\mathbb R^{18},
  \]
  where each $q_k=f''(s_k)$ and $\lambda$ have stress units.

  Positivity is enforced through unconstrained log parameters:
  \[
    \boxed{p_i=\eref\exp(\theta_i)},
    \qquad
    \frac{\partial p}{\partial\theta}=\operatorname{diag}(p).
  \]

  \begin{columns}[T,onlytextwidth]
    \begin{column}{0.47\linewidth}
      \begin{block}{Physical role}
        $\exp(\theta_i)$ is dimensionless; multiplication by $\eref$ restores
        stress units. Choosing $\eref$ near the material scale keeps
        $\theta_i=\log(p_i/\eref)$ numerically moderate.
      \end{block}
    \end{column}
    \begin{column}{0.47\linewidth}
      \begin{block}{Optimization role}
        Every line-search trial satisfies $q_k>0$ and $\lambda>0$ without
        inequality-constrained optimization.
      \end{block}
    \end{column}
  \end{columns}
\end{frame}

\begin{frame}{Shared FEM Derivative Building Block}
  \small
  Let $u\in\mathbb R^{n_u}$ be the nodal displacement and
  \[
    E_h(u,p)\in\mathbb R,
    \qquad
    g(u,p)=\frac{\partial E_h}{\partial u}\in\mathbb R^{n_u}.
  \]
  The common material-sensitivity matrix is
  \[
    \boxed{
      \Eup
      :=\frac{\partial g}{\partial p}
      =\frac{\partial^2E_h}{\partial u\,\partial p}
      \in\mathbb R^{n_u\times18}
    }.
  \]

  Column $i$ is the change in the complete nodal internal-force vector caused
  by a unit change in physical parameter $p_i$, with $u$ held fixed.

  \begin{alertblock}{Why this matters}
    Direct stress fitting and reaction-force fitting use the same $\Eup$.
    The second experiment adds only the response of the equilibrium state to
    a parameter change.
  \end{alertblock}
\end{frame}

\section{Experiment I: Direct Stress Fitting}

\begin{frame}{Constitutive Dataset: Prescribed Deformation Gradients}
  \scriptsize
  Each datum contains a known deformation gradient and target stress,
  $\mathcal D_c=(F_c,P_c^\star)$.

  \begin{table}
    \centering
    \renewcommand{\arraystretch}{1.25}
    \begin{tabular}{llll}
      \toprule
      Family & $F$ & $J$ & Purpose \\
      \midrule
      Uniaxial & $\operatorname{diag}(a,1,1)$ & $a$ & axial and volume coupling \\
      Biaxial & $\operatorname{diag}(a,a,1)$ & $a^2$ & different volume path \\
      Volumetric & $aI$ & $a^3$ & direct volume excitation \\
      Isochoric & $\operatorname{diag}(a,a^{-1},1)$ & $1$ & isolate stretch response \\
      Simple shear & $I+\gamma e_1e_2^T$ & $1$ & off-diagonal response \\
      Rotated distinct & $R_L\operatorname{diag}(e^{\ell_i})R_R^T$ &
        $e^{\sum_i\ell_i}$ & general $F$ and isotropy \\
      Knot-aligned & $\operatorname{diag}(s_k,1,1)$ & $s_k$ & excite every spline basis \\
      \bottomrule
    \end{tabular}
  \end{table}

  \vspace{0.4em}
  \begin{block}{Coverage}
    60 training states and 49 disjoint holdout states cover compression,
    tension, volume change, isochoric distortion, shear, and general rotations.
  \end{block}
\end{frame}

\begin{frame}{Direct Objective and Key Jacobian}
  \small
  For each prescribed $F_c$, fit all nine first-Piola components:
  \[
    r_c(\theta)
    =\frac{\operatorname{vec}\!\left[
      P_{\mathrm{SP}}(F_c;p(\theta))-P_{\mathrm{NH}}(F_c)
    \right]}{\eref},
    \qquad
    \Phi=\frac12\sum_c\norm{r_c}^2.
  \]

  Under affine FEM kinematics, let
  $B=\partial u/\partial\operatorname{vec}F$ and let $V_0$ be the reference
  volume. Virtual work gives
  \[
    \operatorname{vec}P=\frac1{V_0}B^Tg.
  \]
  Since $F_c$ is observed and independent of $p$,
  \[
    \boxed{
      \frac{\partial r_c}{\partial\theta}
      =\frac1{V_0\eref}
      B^T\Eup\operatorname{diag}(p)
    }.
  \]

  \begin{block}{No static solve}
    All nodal displacements are prescribed by $F_c$; there are no unknown
    free DOFs and no state sensitivity $du^\star/dp$.
  \end{block}
\end{frame}

\begin{frame}{Direct-Fit Diagnostics}
  \centering
  \includegraphics[
    width=0.98\linewidth,
    height=0.76\textheight,
    keepaspectratio
  ]{../../systematic_poking_fit_neo_hookean/output/fit_diagnostics.png}

  \vspace{-0.35em}
  \scriptsize
  Fitted curvature follows the Neo-Hookean samples, Gauss--Newton converges
  rapidly, and unseen stress components lie close to the identity line.
\end{frame}

\begin{frame}{Direct-Fit Result and Verification}
  \small
  \begin{table}
    \centering
    \begin{tabular}{lrrr}
      \toprule
      Split & Initial RMSE & Fitted RMSE & Reduction \\
      \midrule
      Training & $1.84\times10^{-1}$ & $3.06\times10^{-4}$ & $599\times$ \\
      Holdout & $1.47\times10^{-1}$ & $2.44\times10^{-4}$ & $603\times$ \\
      \bottomrule
    \end{tabular}
  \end{table}

  \vspace{0.5em}
  \begin{columns}[T,onlytextwidth]
    \begin{column}{0.47\linewidth}
      \textbf{Optimization}
      \begin{itemize}
        \item 6 Gauss--Newton iterations
        \item rank $=18$; condition number $\approx800$
        \item all fitted physical parameters positive
      \end{itemize}
    \end{column}
    \begin{column}{0.47\linewidth}
      \textbf{Independent checks}
      \begin{itemize}
        \item Jacobian FD relative error: $4.8\times10^{-10}$
        \item physical-parameter linear oracle agrees to
          $5.7\times10^{-14}$ relative error
        \item train and holdout errors have the same scale
      \end{itemize}
    \end{column}
  \end{columns}

  \begin{block}{Conclusion of Experiment I}
    The material model, parameter derivatives, stress extraction, and outer
    optimizer are correct before equilibrium is introduced.
  \end{block}
\end{frame}

\section{Experiment II: Static Reaction-Force Fitting}

\begin{frame}{What Changes When Only Boundary Data Are Observed?}
  \small
  Split displacement DOFs into unknown free entries and prescribed entries:
  \[
    u=\begin{bmatrix}u_f\\u_c\end{bmatrix},
    \qquad
    u_f^\star(p)=\arg\min_{u_f}E_h(u_f,u_c;p).
  \]
  The converged state satisfies
  \[
    \boxed{g_f(u_f^\star,u_c,p)=0}.
  \]

  The measured scalar reaction is extracted from the full force vector by
  \[
    \boxed{y(p)=c^Tg(u^\star(p),p)}.
  \]
  Here $c\in\mathbb R^{n_u}$ is a selector: it is $1$ on the measured
  constrained DOFs and $0$ elsewhere. Therefore $c^Tg$ sums, for example,
  all top-face $y$ reactions in an axial test.

  \begin{alertblock}{New dependency}
    Changing $p$ changes both the material force law and the relaxed state
    $u_f^\star(p)$.
  \end{alertblock}
\end{frame}

\begin{frame}{Static Load Cases and Fitting Targets}
  \scriptsize
  The unit cube uses 8 trilinear elements, 27 vertices, and 81 displacement
  DOFs.

  \begin{table}
    \centering
    \renewcommand{\arraystretch}{1.35}
    \begin{tabular}{p{0.18\linewidth}p{0.45\linewidth}p{0.25\linewidth}}
      \toprule
      Protocol & Prescribed boundary motion & Fitted scalar $y$ \\
      \midrule
      Free uniaxial &
        $u_y=0$ at the bottom, $u_y=a-1$ at the top; three scalar pins remove
        tangential rigid modes &
        top-face axial resultant $\sum g_{v,y}$ \\
      Confined uniaxial &
        same axial motion, with $u_x=u_z=0$ throughout &
        top-face axial resultant $\sum g_{v,y}$ \\
      Simple shear &
        bottom clamped; top $u_x=\gamma$ and $u_y=0$ &
        top $x$-resultant $\sum g_{v,x}$ \\
      \bottomrule
    \end{tabular}
  \end{table}

  \begin{block}{Why three protocols?}
    Free uniaxial observes material-dependent lateral relaxation, confined
    uniaxial strongly excites volume change, and shear adds an independent
    distortion path.
  \end{block}
\end{frame}

\begin{frame}{Static Dataset Geometry: Uniaxial Cases}
  \centering
  \includegraphics[
    width=0.93\linewidth,
    height=0.72\textheight,
    keepaspectratio
  ]{../output/uniaxial_static_solve_obj/preview.png}

  \vspace{-0.35em}
  \small
  Target Neo-Hookean static solves at $a=2^{-1/2}\approx0.7071$.
  The free case expands laterally; the confined case retains its
  $1\times1$ cross-section.
\end{frame}

\begin{frame}{Static Dataset Geometry: Simple Shear}
  \centering
  \includegraphics[
    width=0.93\linewidth,
    height=0.72\textheight,
    keepaspectratio
  ]{../output/shear_static_solve_obj/preview.png}

  \vspace{-0.35em}
  \small
  Target Neo-Hookean static solve at $\gamma=0.8$.
  Interior and unconstrained lateral DOFs relax while the top platen motion
  is prescribed.
\end{frame}

\begin{frame}{Dataset Split and Objective}
  \small
  \begin{table}
    \centering
    \begin{tabular}{lrrr}
      \toprule
      Protocol & Train & Validation & Final test \\
      \midrule
      Free uniaxial & 16 & 16 & 32 \\
      Confined uniaxial & 16 & 16 & 32 \\
      Simple shear & 4 & 4 & 8 \\
      \midrule
      Total & 36 & 36 & 72 \\
      \bottomrule
    \end{tabular}
  \end{table}

  Training uses non-rest spline knots; validation uses interval midpoints;
  the sealed final test uses quarter points. The normalized residual is
  \[
    r_j(\theta)=\frac{y_j(p(\theta))-y_j^\star}{\eref}.
  \]
  The fitted objective includes log-curvature smoothness, with $\alpha$
  selected by validation RMSE:
  \[
    \boxed{
      L(\theta)
      =\frac12\sum_{j\in\mathrm{train}}r_j^2
      +\frac{\alpha}{2}\norm{D_2\theta_q}^2
    },
    \qquad \alpha=10^{-3}.
  \]
\end{frame}

\begin{frame}{The Essential Implicit-Differentiation Result}
  \small
  Differentiate free equilibrium with respect to all physical parameters:
  \[
    g_f(u_f^\star,u_c,p)=0
    \quad\Longrightarrow\quad
    \boxed{
      \frac{du_f^\star}{dp}=-\Kff^{-1}\Eupf
    },
  \]
  where
  \[
    \Kff=\frac{\partial g_f}{\partial u_f},
    \qquad
    \Eupf=\frac{\partial g_f}{\partial p}.
  \]

  Apply the chain rule to $y(p)=c^Tg(u^\star(p),p)$:
  \[
    \boxed{
      \frac{dy}{dp}
      =c^T\left(
        \Eup-K_{:f}\Kff^{-1}\Eupf
      \right)
    }.
  \]
  The first term is the direct material effect; the second is the induced
  equilibrium-relaxation effect. Finally,
  \[
    \frac{\partial r}{\partial\theta}
    =\frac1{\eref}\frac{dy}{dp}\operatorname{diag}(p).
  \]
\end{frame}

\begin{frame}{Nested Optimization Pipeline}
  \small
  For every outer evaluation of $\theta$:
  \begin{enumerate}
    \item Convert to positive physical parameters
      $p=\eref\exp(\theta)$.
    \item For every load case, solve
      $g_f(u_f^\star,u_c,p)=0$ with libpgo's Newton solver.
    \item Assemble reactions, $\Kff$, and $\Eup$ at the converged state.
    \item Apply the implicit reaction Jacobian and stack all residual rows.
    \item Append the smoothness residual and take a damped Gauss--Newton step.
  \end{enumerate}

  \[
    \left[J^TJ+\gamma\operatorname{diag}(J^TJ)\right]\Delta\theta
    =-J^Tr.
  \]

  \begin{alertblock}{Important implementation choice}
    The derivative is taken through the converged equilibrium equation, not
    by backpropagating through the sequence of Newton iterations. Every
    line-search trial nevertheless requires fresh static solves.
  \end{alertblock}
\end{frame}

\section{Results, Verification, and Outlook}

\begin{frame}{Reaction-Fit Diagnostics}
  \centering
  \includegraphics[
    width=0.98\linewidth,
    height=0.76\textheight,
    keepaspectratio
  ]{../output/fit_diagnostics.png}

  \vspace{-0.35em}
  \scriptsize
  The recovered curvature tracks the target, the nested objective converges,
  and reaction curves agree across uniaxial and shear protocols.
\end{frame}

\begin{frame}{Quantitative Results Across Both Experiments}
  \scriptsize
  All errors are normalized by $\eref$; the observables differ, so stress and
  reaction RMSE should be read as separate experiment-level metrics.

  \begin{table}
    \centering
    \renewcommand{\arraystretch}{1.2}
    \begin{tabular}{llrr}
      \toprule
      Experiment & Split & Initial RMSE & Fitted RMSE \\
      \midrule
      Direct stress & Training & $1.84\times10^{-1}$ & $3.06\times10^{-4}$ \\
      Direct stress & Holdout & $1.47\times10^{-1}$ & $2.44\times10^{-4}$ \\
      \midrule
      Static reaction & Training & $3.41\times10^{-1}$ & $2.47\times10^{-4}$ \\
      Static reaction & Validation & $3.05\times10^{-1}$ & $2.49\times10^{-4}$ \\
      Static reaction & Final test & $3.06\times10^{-1}$ & $2.35\times10^{-4}$ \\
      \bottomrule
    \end{tabular}
  \end{table}

  \begin{columns}[T,onlytextwidth]
    \begin{column}{0.47\linewidth}
      \begin{block}{Direct stress}
        6 iterations, full rank 18, condition $\approx800$.
      \end{block}
    \end{column}
    \begin{column}{0.47\linewidth}
      \begin{block}{Static reaction}
        7 iterations, full rank 18, condition $\approx458$, and maximum free
        equilibrium residual below $9.7\times10^{-7}$.
      \end{block}
    \end{column}
  \end{columns}
\end{frame}

\begin{frame}{Verification Hierarchy}
  \small
  \begin{table}
    \centering
    \begin{tabular}{lll}
      \toprule
      Check & Direct stress & Static reaction \\
      \midrule
      Analytic vs. central FD Jacobian & $4.8\times10^{-10}$ &
        $3.17\times10^{-10}$ \\
      Data Jacobian rank & $18/18$ & $18/18$ \\
      Held-out coordinates & 49 states & 72 load cases \\
      Positivity & maintained & maintained \\
      Extra independent check & linear oracle & equilibrium residual \\
      \bottomrule
    \end{tabular}
  \end{table}

  \vspace{0.7em}
  \begin{block}{What the static finite difference actually checks}
    Every perturbation performs fresh nonlinear solves. Agreement therefore
    covers constitutive derivatives, FEM assembly, equilibrium sensitivity,
    reaction selection by $c$, and the log-parameter chain rule together.
  \end{block}
\end{frame}

\begin{frame}{Knot Resolution: Accuracy vs. Conditioning}
  \small
  \begin{columns}[T,onlytextwidth]
    \begin{column}{0.48\linewidth}
      \centering
      \textbf{Direct stress holdout}

      \vspace{0.3em}
      \begin{tabular}{rrr}
        \toprule
        Knots & Condition & RMSE \\
        \midrule
        5 & $70$ & $4.67\times10^{-3}$ \\
        9 & $211$ & $1.00\times10^{-3}$ \\
        17 & $800$ & $2.44\times10^{-4}$ \\
        33 & $3800$ & $7.54\times10^{-5}$ \\
        \bottomrule
      \end{tabular}
    \end{column}
    \begin{column}{0.48\linewidth}
      \centering
      \textbf{Static reaction validation}

      \vspace{0.3em}
      \begin{tabular}{rrr}
        \toprule
        Knots & Condition & RMSE \\
        \midrule
        5 & $45$ & $7.45\times10^{-3}$ \\
        9 & $123$ & $1.28\times10^{-3}$ \\
        17 & $461$ & $4.01\times10^{-4}$ \\
        \bottomrule
      \end{tabular}
    \end{column}
  \end{columns}

  \vspace{0.9em}
  \begin{alertblock}{Tradeoff}
    More knots reduce spline approximation error but weaken conditioning.
    Noisy physical data will require smoothness regularization, controlled
    model selection, or a lower-dimensional parameterization.
  \end{alertblock}
\end{frame}

\begin{frame}{What Is Established---and What Is Not}
  \small
  \begin{columns}[T,onlytextwidth]
    \begin{column}{0.47\linewidth}
      \textbf{\good{Established in these experiments}}
      \begin{itemize}
        \item complete 18-parameter optimization chain
        \item positive-domain material parameterization
        \item direct and equilibrium Jacobians verified by FD
        \item full local rank for the selected datasets
        \item transfer to off-knot held-out cases
      \end{itemize}
    \end{column}
    \begin{column}{0.47\linewidth}
      \textbf{\warn{Still outside the evidence}}
      \begin{itemize}
        \item experimental noise and repeated trials
        \item mesh or geometry mismatch
        \item contact-set changes and indenter friction
        \item extrapolation beyond $s\in[0.5,2]$
        \item inversion and collapse behavior
      \end{itemize}
    \end{column}
  \end{columns}

  \vspace{0.6em}
  \begin{block}{Interpretation}
    Confidence is high in the implemented material and equilibrium derivative
    pipeline, but these synthetic same-mesh experiments are not yet evidence
    for robustness on real poking data.
  \end{block}
\end{frame}

\begin{frame}{Takeaways and Next Step}
  \small
  \begin{enumerate}
    \item Direct stress fitting validates the Systematic Poking
      parameterization and its analytic material Jacobian.
    \item Static fitting introduces the implicit state map
      $p\mapsto u_f^\star(p)$ and the correction
      $-K_{:f}\Kff^{-1}\Eupf$.
    \item Both 18-parameter fits reach held-out normalized RMSE near
      $2.4\times10^{-4}$ with full local rank and FD-verified Jacobians.
    \item The current error floor is primarily spline resolution; increasing
      resolution improves accuracy while worsening conditioning.
  \end{enumerate}

  \vspace{0.5em}
  \begin{alertblock}{Next highest-information experiment}
    Replace the prescribed top face with an indenter, fit a force--depth
    curve through contact equilibrium, and repeat the same residual, FD,
    rank, holdout, mesh, and noise checks.
  \end{alertblock}

  \vfill
  \scriptsize
  Reference: Chen, Zhao, and Barbi\v{c},
  \emph{Capturing Animation-Ready Isotropic Materials Using Systematic Poking},
  ACM TOG 42(6), 2023.
\end{frame}

\end{document}
